% =============================================================================
% Kapitel 5: Evaluation
% =============================================================================
\chapter{Evaluation}
\label{ch:evaluation}

\todo{Einleitenden Text verfassen. Was wird evaluiert und mit welchem Ziel?}

% -----------------------------------------------------------------------------
\section{Evaluationsdesign}
\label{sec:evaluationsdesign}

\todo{Das Design der Evaluation beschreiben. Welche Metriken, Kriterien oder Hypothesen werden überprüft?}

% -----------------------------------------------------------------------------
\section{Testumgebung}
\label{sec:testumgebung}

\todo{Die Testumgebung beschreiben (Hardware, Software, Datensätze, Konfigurationen).}

% -----------------------------------------------------------------------------
\section{Ergebnisse}
\label{sec:ergebnisse}

\todo{Die Ergebnisse der Evaluation präsentieren. Tabellen und Grafiken verwenden.}

% Beispiel für eine Tabelle:
% \begin{table}[htbp]
%     \centering
%     \caption{Ergebnisse der Evaluation}
%     \label{tab:ergebnisse}
%     \begin{tabular}{lcc}
%         \toprule
%         \textbf{Metrik} & \textbf{Wert} & \textbf{Einheit} \\
%         \midrule
%         Genauigkeit & 95,3 & \% \\
%         Laufzeit & 1,2 & s \\
%         \bottomrule
%     \end{tabular}
% \end{table}

% Beispiel für eine Abbildung:
% \begin{figure}[htbp]
%     \centering
%     \includegraphics[width=0.8\textwidth]{bilder/beispiel.png}
%     \caption{Beispielabbildung}
%     \label{fig:beispiel}
% \end{figure}

% -----------------------------------------------------------------------------
\section{Zusammenfassung}
\label{sec:evaluation_zusammenfassung}

\todo{Kurze Zusammenfassung der Evaluationsergebnisse.}
